\documentclass[11pt]{article}

\usepackage[margin=1in]{geometry}
\usepackage[T1]{fontenc}
\usepackage[utf8]{inputenc}
\usepackage{amsmath,amssymb}
\usepackage{booktabs}
\usepackage{enumitem}
\usepackage{graphicx}
\usepackage[hidelinks]{hyperref}
\usepackage{listings}
\usepackage{xcolor}

\newcommand{\CourseName}{CS 498: Machine Perception}
\newcommand{\Term}{Fall 2026}
\newcommand{\HomeworkNumber}{HW0}
\newcommand{\HomeworkTitle}{Perspective, Homography, and Stereo}
\newcommand{\CourseWeight}{10\%}
\newcommand{\ReleaseDate}{September 2, 2026}
\newcommand{\DueDate}{September 11, 2026 at 11:59 p.m. America/Chicago}
\newcommand{\task}[3]{\paragraph{Task #1: #2 [#3 points]}}

\definecolor{codegray}{RGB}{245,245,245}
\lstset{basicstyle=\ttfamily\small,backgroundcolor=\color{codegray},breaklines=true,frame=single,showstringspaces=false}

\title{\CourseName\\\HomeworkNumber: \HomeworkTitle}
\author{\Term}
\date{}

\begin{document}
\maketitle

\begin{center}
\begin{tabular}{ll}
\toprule
Course weight & \CourseWeight \\
Release & \ReleaseDate \\
Due & \DueDate \\
Work mode & Individual \\
\bottomrule
\end{tabular}
\end{center}

\section{Overview}
This homework connects planar projective geometry, full camera projection, and
dense stereo. Students first reason about a planar tennis court, estimate and
apply a homography, recover a camera projection matrix from 2D--3D
correspondences, and project a simple mesh into the scene. They then rectify a
stereo pair and implement window-based disparity estimation.

\paragraph{Dataset note.}
The final tennis-court image, landmarks, dimensions, calibration, and insertion
target will be included in the released dataset bundle.

\section{Learning Objectives}
Students will be able to:
\begin{enumerate}[label=\arabic*.]
  \item formulate homogeneous-coordinate constraints for planar and camera projection;
  \item estimate homography and projection matrices with direct linear transforms;
  \item perform inverse image warping and alpha compositing;
  \item explain rectification, disparity, depth, and the effect of intrinsic scaling;
  \item implement and evaluate sum-of-squared-differences stereo matching.
\end{enumerate}

\section{Setup}
From the \texttt{hw0/} directory, create an isolated environment and run:
\begin{lstlisting}[language=bash]
python -m pip install -e ".[dev]"
cs498-hw0 check
pytest
\end{lstlisting}
The starter API is in \texttt{src/cs498\_hw0/}. Keep its public function
signatures unchanged so the tests and grading tools can import your work.

\section{Tasks}
\task{0}{Environment and data check}{0}
Run the environment check, confirm the tennis-court assets are present, and
record the supplied coordinate conventions.

\task{1}{Court coordinates and homography}{15}
Define the tennis-court plane coordinates and implement normalized DLT to
estimate the homography from court coordinates to image coordinates.

\task{2}{Inverse warping and compositing}{20}
Implement inverse image warping and alpha compositing. Place the supplied
graphic or texture at the instructed metric location on the court.

\task{3}{Camera projection matrix}{20}
Use the supplied 2D--3D landmarks to estimate a $3\times4$ camera projection
matrix. Report reprojection error and discuss degeneracy.

\task{4}{3D object insertion}{15}
Transform a supplied mesh into the tennis-court world frame, project its
vertices, and visualize the result. Clearly state the world-frame convention.

\task{5}{Stereo rectification and intrinsic scaling}{5}
Scale the camera intrinsics consistently with image resizing, compute the
relative camera pose, and rectify the stereo pair using the supplied helper.

\task{6}{SSD disparity}{20}
Implement window-based sum-of-squared-differences matching without calling an
off-the-shelf stereo matcher. Compare with a library baseline.

\task{7}{Analysis}{5}
Discuss failure modes near occlusions, repetitive texture, image boundaries,
and low-texture regions. Explain at least two possible improvements.

\section{Deliverables}
Submit completed starter modules, the requested generated arrays and figures,
and a PDF report containing derivations, visualizations, metrics, and analysis.

\section{Rubric}
\begin{center}
\begin{tabular}{lr}
\toprule
Component & Points \\
\midrule
Homography and compositing & 35 \\
Projection and object insertion & 35 \\
Stereo and analysis & 30 \\
\midrule
Total & 100 \\
\bottomrule
\end{tabular}
\end{center}

\section{Submission}
The Canvas submission shell and exact file manifest will be added before
release. The deadline is \DueDate.

\end{document}
